\scanpage{173}\mspage{28}
plus fines de $X$. On peut alors, pour une telle partition \uline{fixée} \inserted{$\unsure{\mathcal P|P}$}, se proposer \unsure{une|d'une} étude approfondie \inserted{de la cat.} des faisceaux (ou complexes de faisceaux, ou champs) qui sont ``$\mathcal P$-constructibles'' (ou plus généralement, qui sont ``loc. constants'' sur chacun des $X_i$). Ces catégories n'auront des \deleted{Pour préciser di} \inserted{structures} vraiment satisfaisantes que si elles sont stables par les opérations courantes -- telles que $\mathrm R\underline{\mathrm{Hom}}$, ou $\mathrm Rj_*j^*$ où $j:X_i\mathrel{\unsure{-|\longrightarrow}}X$ est une ``cellule'' de la partition, \inserted{etc}. En fait, si $X$ est excellent et qu'on dispose de la résolution des singularités, on sait que les faisceaux \inserted{de torsion} constructibles . (sous réserve d'être premiers aux car résiduelles) sont stables par \deleted{toutes} ces opérations -- mais non avec une partition de $X$ fixée une fois pour \unsure{toute|toutes}. Pour avoir une telle stabilité plus fine, il faut faire des hypothèses très strictes \deleted{sur} ``d'équisingularité'' de la stratification donnée de $X$, le long de ses strates. Je pense néanmoins qu'un \unsure{raffinement|affinement} des techniques connues montrera que $X$ admet des stratifications arbitrairement fines ayant ces propriétés d'équisingularité (et avec des $X_i$ réguliers et connexes, mais peu importe \unsure{pour} notre propos \unsure{présent}\ldots). À titre d'exemple
% archive sheet number: 173
